分解是一种在保持数值不变的同时，以较小的底数（基数）表示一个数的数学技术。本节介绍数的分解与多项式分解。

\subsection{数的分解}
\label{subsec:number-decomp}
固定模数 $q\ge 2$，并记 $\mathbb{Z}_q=\mathbb{Z}/q\mathbb{Z}$。设 $\gamma\in\mathbb{Z}_q$。数的分解将 $\gamma$ 表示为以 $\beta$ 为基数的多个数之和：

$ $

$\gamma = \gamma_1 \dfrac{q}{\beta^1} + \gamma_2 \dfrac{q}{\beta^2} + \cdots + \gamma_\ell \dfrac{q}{\beta^\ell}  $

$ $

\noindent 其中，$\beta\ge 2$ 和 $\ell\ge1$ 均为整数。若要求每个规范代表元 $0\le\gamma<q$ 都具有满足 $\gamma_i\in\{0,1,\ldots,\beta-1\}$ 的唯一展开，应假设 $q=\beta^\ell$。更一般地，$\beta^\ell\mid q$ 仅保证权重为整数；若 $d=q/\beta^\ell>1$，则精确展开存在当且仅当 $d\mid\gamma$。图~\ref{fig:decomp} 展示完整分解；截断情形见 \autoref{subsec:approx-decomp}。
当 $q=\beta^\ell$ 时，每个 $\gamma_i$ 是 $\gamma$ 的基 $\beta$ 展开的一个数位，其中 $i=1$ 为最高位。当 $q$ 和 $\beta$ 均为 $2$ 的幂时，才可用二进制位移解释这些权重。

\begin{figure}[h!]
    \centering
  \includegraphics[width=0.8\linewidth]{figures/decomp.pdf}
  \caption{数的分解示意图。}
  \label{fig:decomp}
\end{figure}

将数 $\gamma$ 按基数 $\beta$、层数 $\ell$ 进行分解，定义如下：


$ $

$\textsf{Decomp}^{\beta, \ell}(\gamma) = (\gamma_1, \gamma_2, \text{ } \cdots , \text{ } \gamma_\ell)$.
 
$ $

数的分解也称为基数分解，其中的底数 $\beta$ 称为基数。


\subsubsection{示例}

假设在 $\mathbb{Z}_{16}$ 中取 $\gamma=13$，希望以基数 $\beta = 2$、层数 $\ell = 4$ 分解 13。其分解如下：

$ $

$13 = 1 \cdot \dfrac{16}{2^1} + 1 \cdot \dfrac{16}{2^2} + 0 \cdot \dfrac{16}{2^3} + 1 \cdot \dfrac{16}{2^4}$

$ $

$\textsf{Decomp}^{2, 4}(13) = (1, 1, 0, 1)$


\subsection{多项式分解}
\label{subsec:poly-decomp}

现在，假设在多项式环 ${\mathbb{Z}_q[x] / (x^n + 1)}$ 中有一个多项式 $f$。因此，$f$ 的每个系数 $c_i$ 都是模 $q$ 的整数。多项式分解将 $f$ 表示为以 $\beta$ 为基数、$\ell$ 为层数的多个多项式之和：


\begin{tcolorbox}[title={\textbf{\tboxlabel{\ref*{subsec:poly-decomp}} 多项式分解}}]

给定 $f\in \mathbb{Z}_q[x]/(x^n+1)$，选取满足 $q=\beta^\ell$ 的整数 $\beta\ge2$、$\ell\ge1$，写为

$ $

$
f=\sum\limits_{i=1}^{\ell} f_i\,\dfrac{q}{\beta^i}, \qquad f_i\in \mathbb{Z}_q[x]/(x^n+1)
$

$ $

其中，每个 $f_i$ 都通过将 $f$ 的每个系数按基数 $\beta$ 分解而得到。若 $f=\sum\limits_j c_j x^j$ 且 $c_j\in\mathbb{Z}_q$，则
$c_j=\sum\limits_{i=1}^{\ell} c_{j,i}\,\dfrac{q}{\beta^i}$，其中 $c_{j,i}\in\{0,\ldots,\beta-1\}$，
并且 $f_i=\sum_j c_{j,i} x^j$。
将多项式 $f$ 按基数 $\beta$、层数 $\ell$ 分解，记为：

$ $

$\textsf{Decomp}^{\beta, \ell}(f) = (f_1, f_2, \text{ } \cdots , \text{ } f_\ell)$
 $ $
\end{tcolorbox}




\subsubsection{示例}

假设在多项式环 $\mathbb{Z}_{16}[x] / (x^4 + 1)$ 中有如下多项式：

$ $

$f = 6x^3 + 3x^2 + 14x + 7 \in \mathbb{Z}_{16}[x] / (x^4 + 1)$

$ $

希望按基数 $\beta = 4$、层数 $\ell = 2$ 分解上述多项式。则多项式 $f$ 的各次项系数分别分解如下：

$ $

${\bm{x}^{\bm{3}}}$: $6 = \color{blue}{1 \cdot \dfrac{16}{4^1}} \color{black}+ \color{red}{2 \cdot \dfrac{16}{4^2}}$

${\bm{x}^{\bm{2}}}$: $3 = \color{blue}{0 \cdot \dfrac{16}{4^1}} \color{black}+ \color{red}{3 \cdot \dfrac{16}{4^2}}$

${\bm{x}^{\bm{1}}}$: $14 = \color{blue}{3 \cdot \dfrac{16}{4^1}} \color{black}+ \color{red}{2 \cdot \dfrac{16}{4^2}}$

${\bm{x}^{\bm{0}}}$: $7 = \color{blue}{1 \cdot \dfrac{16}{4^1}} \color{black}+ \color{red}{3 \cdot \dfrac{16}{4^2}}$

$ $

分解得到的多项式如下：

$f = 6x^3 + 3x^2 + 14x + 7 = \color{blue}{(1x^3 + 0x^2 + 3x + 1) \cdot \dfrac{16}{4^1}} \color{black}+ \color{red}{(2x^3 + 3x^2 + 2x + 3) \cdot \dfrac{16}{4^2}} \color{black}$

$ $

$\textsf{Decomp}^{4, 2}(6x^3 + 3x^2 + 14x + 7) = (1x^3 + 0x^2 + 3x + 1, 2x^3 + 3x^2 + 2x + 3)$

\subsubsection{讨论}

分解之后，多项式的原始系数被拆成了更小的数。这一性质在全同态加密密文的多项式乘法中非常有用，可以降低噪声的增长速度。通常，密文的多项式系数较大，因为它们是均匀随机选取的数。减小这些较大的常数对于控制同态乘法中的噪声增长很重要。\autoref{subsec:tfhe-mult-cipher} 将进一步讨论这一点。


\subsection{近似分解}
\label{subsec:approx-decomp}

\begin{figure}[h!]
    \centering
  \includegraphics[width=0.7\linewidth]{figures/decomp3.pdf}
  \caption{近似分解示意图。}
  \label{fig:decomp3}
\end{figure}

选取满足 $\beta^\ell\mid q$ 的整数 $\beta\ge2$、$\ell\ge1$，并令 $d=q/\beta^\ell$。即使整除条件成立，只要 $d>1$ 且输入不是 $d$ 的倍数，就仍须截断。将 $\gamma$ 的规范代表元舍入到最接近的 $d$ 的整数倍，再将该整数倍模 $q$ 约减，并按上述权重分解。固定最近整数舍入的平局规则后，结果满足
$$
\gamma\equiv\sum_{i=1}^{\ell}\gamma_i\frac{q}{\beta^i}+\varepsilon\pmod q,
\qquad \gamma_i\in\{0,\ldots,\beta-1\},\qquad |\varepsilon|\le\frac{q}{2\beta^\ell}.
$$
残差取圆周舍入误差的一个整数提升。这里必须使用同余：例如 $q=16$、$\beta=2$、$\ell=2$ 时，输入 $15$ 舍入到 $16\equiv0\pmod{16}$，此时 $\varepsilon=-1$。多项式情形逐系数处理。若 $\beta^\ell\nmid q$，上述有理数权重不能定义这里的整数环 Gadget；可以改用 \autoref{subsec:gadget-decomposition} 中的整数 Gadget。


\subsection{Gadget 分解}
\label{subsec:gadget-decomposition}

Gadget 权重和允许的数位集合都是定义的一部分；任意权重不一定能表示每个剩余类。对于任意 $q\ge2$，一种精确的整数选择是 $g_i=\beta^{i-1}$，取 $\ell=\lceil\log_\beta q\rceil$，并使用规范代表元的普通进制数位。它在整数中满足 $\gamma=\sum_{i=1}^{\ell}\gamma_i g_i$。这与下面的降序权重不同，两种权重的数位向量不可混用。对于截断 Gadget，应使用带残差的模 $q$ 重构式，并在后续噪声分析中计入该残差。

Gadget 分解是数的分解（\autoref{subsec:number-decomp}）的一般化形式。在数的分解中，数 $\gamma$ 被分解如下：

$\gamma = \gamma_1 \dfrac{q}{\beta^1} + \gamma_2 \dfrac{q}{\beta^2} + \cdots + \gamma_\ell \dfrac{q}{\beta^\ell} $

$ $

在 Gadget 分解中，按如下方式分解 $\gamma$：

$\gamma = \gamma_1 g_1 + \gamma_2 g_2 + \cdots + \gamma_\ell g_\ell $

$ $

将 $\vec{g} = (g_1, g_2, \cdots, g_\ell)$ 称为 Gadget 向量，并记 $\textsf{Decomp}^{\vec{g}}(\gamma) = (\gamma_1, \gamma_2, \text{ } \cdots , \text{ } \gamma_\ell) $。

$ $

于是，$\gamma = \langle \textsf{Decomp}^{\vec{g}}(\gamma), \vec{g} \rangle $。

$ $

对于数的分解（\autoref{subsec:number-decomp}），其 Gadget 向量为 $\vec{g} = \Bigg(\dfrac{q}{\beta}, \dfrac{q}{\beta^2}, \cdots, \dfrac{q}{\beta^\ell}\Bigg)$。